\documentclass[a4paper]{article}

% ==============================================================================
% 1. ĐỘNG CƠ BIÊN DỊCH VÀ CẤU HÌNH FONT TIMES NEW ROMAN
% ==============================================================================
\usepackage{iftex}

\ifXeTeX
  % Biên dịch bằng XeLaTeX: Dùng font Times New Roman hệ thống
  \usepackage{fontspec}
  \setmainfont{Times New Roman}
  \usepackage[vietnamese]{babel}
\else
  % Biên dịch bằng pdfLaTeX: Dùng chuẩn T5 và gói font Times chuẩn
  \usepackage[utf8]{inputenc}
  \usepackage[vietnamese]{babel}
  \usepackage[T5]{fontenc}
  \usepackage{times} % Font Times New Roman chuẩn
\fi

% ==============================================================================
% 2. CỠ CHỮ 13PT & CĂN LỀ CHUẨN ĐHQG-HCM / ĐHBK
% ==============================================================================
% Thiết lập cỡ chữ cơ sở chính xác là 13pt
\usepackage{scrextend}
\changefontsizes{13pt}

% Căn lề chuẩn văn bản học thuật (Nghị định 30 & Quy chuẩn Luận án/Tiểu luận ĐHBK):
% - Trái (gáy đóng tập): 3.0 cm
% - Phải: 2.0 cm
% - Trên: 2.0 cm
% - Dưới: 2.0 cm
\usepackage{geometry}
\geometry{
  a4paper,
  top=20mm,
  bottom=20mm,
  left=30mm,
  right=20mm,
  headheight=16pt,
  footskip=12mm
}

% ==============================================================================
% 3. GIÃN DÒNG 1.5 LINES & THỤT ĐẦU DÒNG
% ==============================================================================
\usepackage{setspace}
\setstretch{1.5} % Giãn dòng 1.5 lines chuẩn
\setlength{\parindent}{1.0cm} % Thụt đầu dòng 1.0 cm (hoặc 1.27 cm / 0.5 inch)
\setlength{\parskip}{6pt}      % Khoảng cách giữa các đoạn văn

% ==============================================================================
% 4. CÁC GÓI BỔ TRỢ HỌC THUẬT CẦN THIẾT
% ==============================================================================
\usepackage{amsmath, amssymb, amsfonts, mathtools} % Toán học nâng cao
\usepackage{graphicx}                              % Chèn ảnh
\usepackage{booktabs, tabularx, multirow, makecell}% Bảng biểu chuyên nghiệp
\usepackage{cite}                                  % Trích dẫn tài liệu
\usepackage{float}                                 % Cố định vị trí hình/bảng [H]
\usepackage[font=small, labelfont=bf]{caption}     % Tùy biến chú thích hình/bảng
\usepackage{subcaption}                            % Ảnh con (subfigures)
\usepackage{enumitem}                              % Tùy biến danh sách bullet
\setlist{nosep, leftmargin=1.5cm}                  % Danh sách gọn gàng

% Khung câu hỏi / đề bài tập (Tự dựng bằng xcolor và lrbox, không phụ thuộc gói ngoài)
\usepackage[table]{xcolor}
\definecolor{boxbg}{RGB}{245, 248, 252}
\definecolor{boxframe}{RGB}{30, 80, 160}

\newsavebox{\hwbox}
\newenvironment{homeworkbox}[1][Yêu cầu bài tập]{%
  \begin{lrbox}{\hwbox}%
    \begin{minipage}{0.92\textwidth}
      \noindent\textbf{\color{boxframe}\large #1}\par
      \vspace{2mm}\hrule height 0.6pt\vspace{3mm}
      \ignorespaces
}{%
    \end{minipage}%
  \end{lrbox}%
  \begin{center}
    \setlength{\fboxsep}{10pt}%
    \setlength{\fboxrule}{1pt}%
    \fcolorbox{boxframe}{boxbg}{\usebox{\hwbox}}%
  \end{center}
}

% ==============================================================================
% 5. HEADER & FOOTER
% ==============================================================================
\usepackage{fancyhdr}
\pagestyle{fancy}
\fancyhf{}
\fancyhead[L]{\small\textit{GK5007 -- PPNCKH Nâng cao}}
\fancyhead[R]{\small\textit{NCS. Nguyễn Minh Hiếu}}
\fancyfoot[C]{\thepage}
\renewcommand{\headrulewidth}{0.4pt}
\renewcommand{\footrulewidth}{0pt}

% ==============================================================================
% 6. HYPERREF (LIÊN KẾT & BOOKMARK)
% ==============================================================================
\usepackage{hyperref}
\hypersetup{
  colorlinks=true,
  linkcolor=black,     % Giữ màu đen cho mục lục để tiện in ấn
  citecolor=blue!80!black,
  urlcolor=blue!80!black,
  pdftitle={Bài tập về nhà - PPNCKH Nâng cao (GK5007)},
  pdfauthor={Nguyễn Minh Hiếu - HCMUT}
}

% ==============================================================================
% THÔNG TIN BÀI TẬP VÀ SINH VIÊN / NCS
% ==============================================================================
\newcommand{\CourseCode}{GK5007}
\newcommand{\CourseName}{PHƯƠNG PHÁP NGHIÊN CỨU KHOA HỌC NÂNG CAO}
\newcommand{\AssignmentTitle}{BÀI TẬP VỀ NHÀ}
\newcommand{\AssignmentSubtitle}{Báo cáo thực hành \& Tiểu luận phương pháp luận nghiên cứu}
\newcommand{\StudentName}{NGUYỄN MINH HIẾU}
\newcommand{\StudentID}{2680050}
\newcommand{\Department}{Khoa Khoa học \& Kỹ thuật Máy tính}
\newcommand{\Instructor}{[Tên Giảng viên phụ trách]}
\newcommand{\Semester}{Học kỳ 1 -- Năm học 2026--2027}
\newcommand{\LocationDate}{TP. Hồ Chí Minh, \today}

% ==============================================================================
% BẮT ĐẦU TÀI LIỆU
% ==============================================================================
\begin{document}
\hypersetup{pageanchor=false}

% ------------------------------------------------------------------------------
% TRANG BÌA CHUẨN BÁCH KHOA - HCMUT
% ------------------------------------------------------------------------------
\begin{titlepage}
  \thispagestyle{empty}
  \begin{center}
    {\large \textbf{ĐẠI HỌC QUỐC GIA TP. HỒ CHÍ MINH}}\\[0.2cm]
    {\large \textbf{TRƯỜNG ĐẠI HỌC BÁCH KHOA}}\\[0.2cm]
    {\normalsize \textbf{\Department}}\\[0.5cm]

    \vspace{0.8cm}
    \includegraphics[width=3.2cm]{figures/HCMCUT-logo.png}\\[1.0cm]

    {\large \textbf{\CourseName}}\\[0.2cm]
    {\normalsize Mã học phần: \textbf{\CourseCode}}\\[1.2cm]

    \rule{\linewidth}{1.2pt}\\[0.4cm]
    {\Large \textbf{\AssignmentTitle}}\\[0.3cm]
    {\large \AssignmentSubtitle}\\[0.2cm]
    \rule{\linewidth}{1.2pt}\\[2.0cm]

    \begin{minipage}[t]{0.85\linewidth}
      \begin{tabular}{ll}
        \textbf{Nghiên cứu sinh / Học viên:} & \textbf{\StudentName} \\
        \textbf{Mã số NCS/HV:}               & \StudentID \\
        \textbf{Học kỳ:}                     & \Semester \\
      \end{tabular}
    \end{minipage}

    \vfill
    {\normalsize \textbf{\LocationDate}}
  \end{center}
\end{titlepage}

% ------------------------------------------------------------------------------
% MỤC LỤC
% ------------------------------------------------------------------------------
\hypersetup{pageanchor=true}
\pagenumbering{roman}
\tableofcontents
\newpage

\pagenumbering{arabic}
\setcounter{page}{1}

% ==============================================================================
% NỘI DUNG CHÍNH CỦA BÀI TẬP (TÁCH THÀNH CÁC FILE RIÊNG ĐỂ DỄ CHỈNH SỬA)
% ==============================================================================
\section{Chương 2}
\label{sec:tong_quan}


\section{Phương pháp luận và mô hình đề xuất}
\label{sec:phuong_phap}

\subsection{Mô hình toán học}
Giả sử hệ thống được mô hình hóa dưới dạng đồ thị có hướng $G = (V, E)$, trong đó $V$ đại diện cho tập các nút tính toán và $E$ là các liên kết truyền thông. Hàm mục tiêu tối ưu hóa được phát biểu như sau:
\begin{equation}
  \min_{\mathbf{w}, \mathbf{p}} \quad \mathcal{L}(\mathbf{w}) + \lambda \sum_{i \in V} \mathcal{C}_i(p_i)
  \label{eq:objective}
\end{equation}
Trong đó:
\begin{itemize}
  \item $\mathbf{w}$ biểu diễn véc-tơ tham số của mô hình học máy.
  \item $p_i$ là mức năng lượng phân bổ cho nút thứ $i$.
  \item $\lambda > 0$ là hệ số điều hòa giữa độ chính xác và chi phí năng lượng.
\end{itemize}

\subsection{Bảng tổng hợp tham số và ký hiệu}
Bảng \ref{tab:notation} tóm tắt các ký hiệu toán học chính được sử dụng trong bài tập này.

\begin{table}[H]
  \centering
  \caption{Bảng ký hiệu và thông số nghiên cứu}
  \label{tab:notation}
  \begin{tabularx}{\textwidth}{l c X}
    \toprule
    \textbf{Ký hiệu} & \textbf{Đơn vị} & \textbf{Ý nghĩa học thuật} \\
    \midrule
    $N$ & -- & Số lượng nút tính toán trong cụm \\
    $B$ & MHz & Băng thông kênh truyền thông vô tuyến \\
    $P_{\max}$ & Watt & Công suất phát tối đa của thiết bị \\
    $T_{\text{delay}}$ & ms & Độ trễ lan truyền tối đa cho phép \\
    $\mathcal{L}(\cdot)$ & -- & Hàm mất mát (Loss function) huấn luyện \\
    \bottomrule
  \end{tabularx}
\end{table}

\section{Kết quả thực nghiệm và thảo luận}
\label{sec:thuc_nghiem}

\subsection{Thiết lập thực nghiệm}
Mô phỏng được xây dựng trên bộ dữ liệu thực tế nhằm đánh giá độ hội tụ và tính ổn định của giải thuật. Kết quả so sánh giữa phương pháp đề xuất và các phương pháp chuẩn (Baseline) được trình bày chi tiết.

\begin{table}[H]
  \centering
  \caption{So sánh hiệu năng giữa giải thuật đề xuất và các Baseline}
  \label{tab:results}
  \begin{tabular}{lcccc}
    \toprule
    \textbf{Phương pháp} & \textbf{Độ chính xác (\%)} & \textbf{Thời gian (s)} & \textbf{Năng lượng (J)} & \textbf{Cải thiện} \\
    \midrule
    Baseline A & 82.4 & 145.2 & 320.5 & Gốc \\
    Baseline B & 86.1 & 112.8 & 275.0 & +4.5\% \\
    \textbf{Đề xuất (Ours)} & \textbf{91.8} & \textbf{78.4} & \textbf{189.2} & \textbf{+11.4\%} \\
    \bottomrule
  \end{tabular}
\end{table}

\subsection{Thảo luận kết quả}
Từ Bảng \ref{tab:results}, ta nhận thấy phương pháp đề xuất giảm thời gian xử lý tới $46.0\%$ so với Baseline A trong khi vẫn nâng cao độ chính xác lên $91.8\%$. Điều này minh chứng tính đúng đắn của giả thuyết nghiên cứu đã đặt ra ban đầu.

\section{Kết luận và hướng phát triển}
\label{sec:ket_luan}

Bài tập này đã giải quyết thành công các yêu cầu cốt lõi về phương pháp luận nghiên cứu khoa học. Trong các bước tiếp theo, mô hình sẽ được mở rộng trên phần cứng thực tế và kiểm thử với các kịch bản nhiễu phức tạp.


% ==============================================================================
% TÀI LIỆU THAM KHẢO
% ==============================================================================
\newpage
\bibliographystyle{ieeetr}
\bibliography{references}

\end{document}
